\documentclass[runningheads]{llncs}
\usepackage[T1]{fontenc}
\usepackage{amsmath,amssymb}
\usepackage{booktabs}
\usepackage{graphicx}

\begin{document}

\title{Supplementary Material:\\Privacy-Preserving Federated Learning for Imbalanced Medical Applications via Distribution-Aware Aggregation}

\author{Anonymized Authors}
\authorrunning{Anonymized Author et al.}
\institute{Anonymized Affiliations \\
    \email{email@anonymized.com}}

\maketitle

\appendix

\section{Convergence Analysis}\label{app:convergence}

This appendix provides the detailed theoretical justification for the convergence remark in the main paper (Section~4.2).
We argue that each phase of the proposed framework converges independently, that the mid-training switch at round~$R^*$ is well-founded, and that the inverse-frequency weighting targets a principled objective.

\subsection{Notation and assumptions}

We adopt the standard federated optimisation setting.
Let $F_i(\theta)$ be client~$i$'s local objective and
$F(\theta)=\sum_{i=1}^N p_i\,F_i(\theta)$
the global objective with aggregation weights~$p_i$.
We assume:
\begin{enumerate}
\item[\textbf{A1.}] \emph{$L$-smoothness.} Each $F_i$ is $L$-smooth: $\|\nabla F_i(\theta)-\nabla F_i(\theta')\|\le L\|\theta-\theta'\|$ for all~$\theta,\theta'$.
\item[\textbf{A2.}] \emph{Bounded variance.} The stochastic gradient $g_i(\theta)$ satisfies $\mathbb{E}\|g_i(\theta)-\nabla F_i(\theta)\|^2\le\sigma^2$.
\item[\textbf{A3.}] \emph{Bounded heterogeneity.} There exists $\zeta^2\ge0$ such that $\sum_i p_i\|\nabla F_i(\theta)-\nabla F(\theta)\|^2\le\zeta^2$ for all~$\theta$.
\end{enumerate}

\subsection{Phase-wise convergence via Li et al.\ and FedNova}\label{app:phase_convergence}

\textbf{Phase~1 (rounds $1$ to $R^*$)} runs standard FedAvg with dataset-size weights $p_i^{(1)}=n_i/\sum_j n_j$.
\textbf{Phase~2 (rounds $R^*\!+\!1$ to $T$)} runs weighted FedAvg with inverse-frequency weights $p_i^{(2)}=w_i$ as defined in Eq.~(5) of the main paper.

In both phases the aggregation weights are \emph{fixed} across rounds (Phase~2 weights are computed once from the histogram estimate and never updated).
Each phase is therefore a standard instance of weighted FedAvg with a well-defined global objective $F^{(s)}(\theta)=\sum_i p_i^{(s)} F_i(\theta)$, $s\in\{1,2\}$.

\paragraph{Li et al.~\cite{ref_li_convergence} (ICLR 2020).}
This work provides the first rigorous convergence proof of FedAvg on non-IID data under partial client participation.
Their Theorem~2 (full participation) and Theorem~3 (partial participation) establish that, under assumptions A1--A3, FedAvg with learning rate $\eta_r=O(1/r)$ and arbitrary sampling probabilities $\{p_i\}$ achieves:
\[
\frac{1}{T}\sum_{r=1}^{T}\mathbb{E}\|\nabla F(\theta^r)\|^2
\;\le\;
O\!\left(\frac{1}{T}\right),
\]
where the constant depends on $L$, $\sigma^2$, $\zeta^2$, the number of local steps~$E$, and the non-uniformity of the sampling weights (quantified by $\nu=N\!\cdot\!\max_i p_i$ and $\varsigma=N\!\cdot\!\min_i p_i$).

The key insight for our framework is that \emph{the convergence guarantee holds for any fixed weight vector} $\{p_i\}$, not just the dataset-size-proportional one.
In Phase~1, setting $p_i=n_i/\sum_j n_j$ recovers the original FedAvg analysis.
In Phase~2, setting $p_i=w_i$ (the inverse-frequency weights) produces a different global objective $F^{(2)}$ whose minimiser better reflects class-balanced performance; the same convergence rate $O(1/T)$ applies, albeit with different constants reflecting the new weight non-uniformity.

\paragraph{Wang et al.~(FedNova)~\cite{ref_fednova} (NeurIPS 2020).}
FedNova introduces a complementary perspective: when clients perform heterogeneous numbers of local steps~$\tau_i$, naive weighted averaging converges to a stationary point of a \emph{mismatched} objective $\tilde{F}\neq F$ (their Theorem~1).
The mismatch $\|F-\tilde{F}\|$ is explicitly characterised as a function of the weights $\{p_i\}$ and local step counts $\{\tau_i\}$.
FedNova eliminates this inconsistency by normalising the weights: $p_i^{\text{nova}} \propto p_i/\tau_i$.

This result is relevant to our framework for two reasons:
\begin{enumerate}
\item It shows that \emph{changing the aggregation weights changes the effective objective} being optimised.
Our Phase~1$\to$Phase~2 switch deliberately changes the objective from a dataset-size-biased one to a class-balanced one: this is a feature, not a bug.

\item It provides convergence guarantees for weighted FedAvg with arbitrary weight vectors under heterogeneous local computation.
Even if clients perform different numbers of local steps in practice (e.g.\ due to hardware heterogeneity), the FedNova normalisation can be applied on top of our inverse-frequency weights, preserving convergence.
\end{enumerate}

\subsection{Heterogeneous clients and stochastic weights via Fraboni et al.}\label{app:fraboni}

\paragraph{Fraboni et al.~\cite{ref_fraboni2023} (JMLR 2023).}
This work provides the most general convergence framework for federated optimisation that is directly applicable to our setting.
Their analysis covers:
\begin{itemize}
\item \emph{Stochastic aggregation weights}: the weights $\{p_i^r\}$ may vary across rounds (e.g.\ due to random client sampling), and the convergence analysis accounts for the variance of the weights themselves.
\item \emph{Heterogeneous client updates}: clients may perform different amounts of local computation, use different local optimisers, and update asynchronously.
\item \emph{Arbitrary weight distributions}: the theory does not require weights to be proportional to dataset size; any weight scheme with bounded second moments is covered.
\end{itemize}

Their main theorem (Theorem~1) establishes convergence to a neighbourhood of the weighted objective's minimiser at rate $O(1/\sqrt{T})$ for non-convex objectives, with the neighbourhood radius depending on the variance of the stochastic weights and the heterogeneity of client updates.

This result is particularly relevant to our framework because:
\begin{enumerate}
\item \textbf{Phase~2 with partial participation.}
In Phase~2, the effective weight of client~$i$ at round~$r$ is $w_i^r$ if $P_i\in\mathcal{P}^r$ and $0$ otherwise.
The random client sampling makes the weights stochastic even though the inverse-frequency parameters $\omega_\ell$ are fixed.
Fraboni et al.'s framework handles exactly this situation: the expected weight $\mathbb{E}[w_i^r]$ defines the target objective, and the variance of $w_i^r$ (due to sampling) contributes to the convergence neighbourhood.

\item \textbf{DP-induced weight perturbation.}
The inverse-frequency weights $\omega_\ell$ are derived from a noisy histogram estimate $\hat{\mathbf{y}}$.
The DP noise perturbs $\omega_\ell$ away from the true inverse-frequency values, introducing an additional source of weight randomness.
As argued in the main paper, the mapping $\hat{p}\mapsto\omega$ is Lipschitz-continuous for $\hat{p}_\ell$ bounded away from zero (ensured by the floor $\delta>0$ applied before normalisation; see Eq.~(3) of the main paper), so the weight perturbation is bounded:
\[
\|\omega(\hat{p}) - \omega(p^*)\|
\;\le\;
L_\omega \cdot \|\hat{p}-p^*\|
\;\le\;
L_\omega \cdot O\!\bigl(\sqrt{\mathrm{Var}(\hat{p})}\bigr),
\]
where $L_\omega$ is the Lipschitz constant of the inverse-frequency mapping and $\mathrm{Var}(\hat{p})$ is the estimation variance decomposed in Eq.~(6) of the main paper.
Under Fraboni et al.'s framework, this bounded perturbation translates to a proportionally bounded enlargement of the convergence neighbourhood, meaning that convergence is preserved up to an additive error controlled by the histogram estimation quality.

\item \textbf{The switch as a special case.}
The Phase~1$\to$Phase~2 transition can be viewed as a single run with time-varying weights: $p_i^r = p_i^{(1)}$ for $r\le R^*$ and $p_i^r = w_i^r$ for $r>R^*$.
Since Fraboni et al.'s theory allows time-varying stochastic weights with bounded variance, the entire two-phase protocol (including the switch) falls within their convergence guarantees.
The fact that the switch is a one-time event means the weight variance is actually \emph{lower} than in settings with continuously varying weights.
\end{enumerate}

\subsection{Inverse-frequency weights and the class-balanced loss}\label{app:cui}

\paragraph{Cui et al.~\cite{ref_cui2019} (CVPR 2019).}
This work provides the theoretical motivation for \emph{why} inverse-frequency weighting is the right choice for class-imbalanced learning, independent of the federated setting.

\paragraph{Effective number of samples.}
Cui et al.\ introduce the concept of the \emph{effective number of samples} for a class with $n$ training examples:
\[
E_n = \frac{1-\beta^n}{1-\beta},
\qquad \beta\in[0,1).
\]
The key insight is that as more data is collected for a given class, the marginal benefit of each additional sample diminishes due to overlap in the data manifold: new samples are increasingly likely to be near-duplicates of existing ones.
The effective number $E_n$ captures this saturation effect and provides a principled re-weighting factor.

\paragraph{Class-balanced loss.}
The class-balanced loss assigns each class~$\ell$ a weight inversely proportional to its effective number of samples:
\[
\mathcal{L}_{\text{CB}} = \frac{1}{E_{n_\ell}} \cdot \mathcal{L}(y, f(\mathbf{x};\theta)).
\]
In the limit $\beta\to 1$, the effective number reduces to $E_n\approx n$, and the class-balanced loss becomes equivalent to inverse-frequency weighting: $\text{weight}_\ell\propto 1/n_\ell$.
This is exactly the weighting strategy adopted in our Phase~2 aggregation.

\paragraph{Connection to our framework.}
In our setting, $\hat{p}_\ell = \hat{y}_\ell / \sum_{\ell'}\hat{y}_{\ell'}$ estimates the global class frequency from the DP histogram, and the aggregation weight $\omega_\ell\propto 1/\hat{p}_\ell$ (Eq.~(3) of the main paper) implements inverse-frequency weighting at the class level.
When these class-level weights are combined with client label counts to produce client-level weights $w_i^r$ (Eq.~(5)), the effective global objective being optimised in Phase~2 becomes:
\[
F^{(2)}(\theta) = \sum_{i=1}^{N} w_i \, F_i(\theta)
= \sum_{i=1}^{N} \frac{\sum_\ell \omega_\ell\,c_{i,\ell}}{\sum_j\sum_\ell \omega_\ell\,c_{j,\ell}} \, F_i(\theta).
\]
Since $\omega_\ell\propto 1/\hat{p}_\ell$, clients whose data is concentrated on rare classes receive higher weight, and the per-class contribution to the global gradient becomes approximately uniform across classes.
This is precisely the class-balanced objective of Cui et al., extended from the centralised to the federated setting.

The practical consequence is that Phase~2 does not merely ``rebalance'' in an ad hoc manner: it optimises a well-motivated objective that has been shown~\cite{ref_cui2019} to yield significant improvements on long-tailed recognition benchmarks (up to $+4\%$ on CIFAR-100-LT and $+2.5\%$ on ImageNet-LT).
In the medical context, where class imbalance is endemic (e.g.\ rare arrhythmia subtypes, uncommon skin lesion categories), this principled re-weighting is especially valuable.

\subsection{Summary of the convergence argument}

Combining the results above, the convergence of our two-phase framework rests on three pillars:

\begin{enumerate}
\item \textbf{Per-phase convergence (Li et al.\ + FedNova):}
Each phase independently optimises a weighted federated objective with fixed weights and inherits the $O(1/T)$ convergence rate of weighted FedAvg under standard smoothness and variance assumptions.

\item \textbf{Robustness to weight perturbation and stochasticity (Fraboni et al.):}
The stochastic nature of the Phase~2 weights (due to client sampling and DP noise) is covered by the general theory for federated optimisation with stochastic aggregation weights.
Convergence is preserved up to a neighbourhood whose radius is controlled by the histogram estimation variance.

\item \textbf{Principled weighting objective (Cui et al.):}
The inverse-frequency weights used in Phase~2 are not arbitrary: they make the effective global objective approximate the class-balanced loss, a well-established re-weighting strategy for imbalanced learning that treats all classes equally regardless of prevalence.
\end{enumerate}

The switch at round $R^*$ constitutes a warm start: the model continues from $\theta^{R^*}$, which provides a better initialisation than random weights since Phase~1 has already learned useful feature representations from the data.
Formally, the entire two-phase protocol is covered by Fraboni et al.'s framework as a single run with time-varying weights (Section~\ref{app:fraboni}), so convergence of Phase~2 is guaranteed regardless of the initialisation quality; the warm start simply offers a practical speed-up.
Since the switch is a one-time event and the Phase~2 weights are fixed thereafter, the convergence theory applies from round $R^*\!+\!1$ onward.

%
% ---- Bibliography ----
%
\begin{thebibliography}{5}

\bibitem{ref_li_convergence}
Li, X., Huang, K., Yang, W., Wang, S., Zhang, Z.: On the convergence of FedAvg on non-IID data. In: ICLR (2020)

\bibitem{ref_fednova}
Wang, J., Liu, Q., Liang, H., Joshi, G., Poor, H.V.: Tackling the objective inconsistency problem in heterogeneous federated optimization. In: NeurIPS, vol.~33, pp.~7611--7623 (2020)

\bibitem{ref_fraboni2023}
Fraboni, Y., Vidal, R., Kameni, L., Lorenzi, M.: A general theory for federated optimization with asynchronous and heterogeneous clients updates. J. Mach. Learn. Res. \textbf{24}(110), 1--48 (2023)

\bibitem{ref_cui2019}
Cui, Y., Jia, M., Lin, T.-Y., Song, Y., Belongie, S.: Class-balanced loss based on effective number of samples. In: CVPR, pp.~9268--9277 (2019)

\end{thebibliography}

\end{document}
